% Rúbrica: Modelación supervisada (selección del modelo y evaluación de su desempeño). Fuente:
% notebooks/01_modelo_riesgo.ipynb, §5 a §8; 00, §3.2, §3.4 y §3.8; backend/modelado/, backend/api/scoring.py;
% docs/evidencia/ (simulacion_123.txt, particion-alternativa.json, prueba-externa.json,
% bootstrap-prueba-externa.json y metacritic-por-banda.md).
\section{Modelación y evaluación}\label{sec:modelacion}

\subsection{El modelo y sus conjuntos de variables}\label{sec:modelacion-conjuntos}

El modelo es una regresión logística\footnote{Un modelo lineal de clasificación: le da un peso a
cada variable, las suma y convierte el resultado en un score entre 0 y 1.} con las variables
estandarizadas y pesos de clase balanceados\footnote{Con \texttt{class\_weight="balanced"}, las
pocas reseñas con $Y = 1$ pesan en el ajuste tanto como todas las demás juntas, para que la clase
rara no se pierda.}, sin ajuste de hiperparámetros. Se eligió por tres razones. Usa pocas
variables, cinco y todas del juego. Sus coeficientes se leen directo, y son los que alimentan la
ficha: cada factor que muestra el sitio es un coeficiente por el valor del juego
(sección~\ref{sec:interpretabilidad}). Y aunque hay \cifraResenasEntrenamiento{} reseñas, lo que
varía son \cifraJuegosEntrenamiento{} juegos, y con tan pocos un modelo más flexible tendría poco
de dónde aprender sin memorizarlos. La comparación contra otros algoritmos viene de la Parte A, con
los modelos de texto (sección~\ref{sec:modelacion-texto}). La misma función,
\texttt{construir\_pipeline}, arma el modelo que sirve la API y el de los notebooks, y se probó con
tres conjuntos de variables, siempre contra el clasificador trivial (tabla~\ref{tab:conjuntos}).

\begin{table}[htbp]
  \centering\small
  % Generado por documento/generar_figuras.py. No se edita a mano.
\begin{tabularx}{\textwidth}{l >{\raggedright\arraybackslash}X rrrr}
\toprule
Conjunto & Variables & Cuántas & PR-AUC & \begin{tabular}[b]{@{}r@{}}Desv. entre\\folds\end{tabular} & \begin{tabular}[b]{@{}r@{}}Veces el\\trivial\end{tabular} \\
\midrule
Trivial & ninguna & 0 & 0.0219 & 0.0037 & 1.0× \\
Juego (el del sitio) & gratuidad, precio, descuento y crítica & 5 & 0.0694 & 0.0415 & 3.2× \\
Compra & las de juego y los juegos del autor & 6 & 0.0709 & 0.0270 & 3.2× \\
Completo & las de compra, la bandera de perfil privado y las posteriores a la reseña & 11 & 0.0849 & 0.0207 & 3.9× \\
\bottomrule
\end{tabularx}

  \caption{PR-AUC del mismo modelo con cada conjunto de variables: media y desviación estándar de
  los \cifraPliegues{} folds de la partición congelada. Fuente: \texttt{evaluar\_gkf} de
  \texttt{backend/modelado/entrenar\_baseline.py} sobre data-v1.}
  \label{tab:conjuntos}
\end{table}

El conjunto completo da el PR-AUC más alto, \cifraPRAUCCompleto{}, pero no se puede usar: sus
variables posteriores a la reseña no están disponibles al momento de la compra
(sección~\ref{sec:procesamiento-variables}). La comparación que importa es entre juego y compra.

\subsection{La decisión B+: no forzar una personalización que los datos no justifican}\label{sec:modelacion-bmas}

El proyecto se quedó con el conjunto juego, la variante que llamó B+, por dos razones. La primera
es estadística. Compra le agrega al conjunto juego la biblioteca del autor, lo único del jugador
que el sitio podría pedir en un formulario, y con eso el PR-AUC sube \cifraDiferenciaCompraJuego{},
muy por debajo de la desviación entre folds del modelo (\cifraPRAUCModeloStd{}). Compra le gana a
juego en \cifraFoldsCompraGana{} de los \cifraPliegues{} folds, y en \cifraParticionesCompraGana{}
de \cifraParticionesSimuladas{} particiones distintas de los juegos: más o menos la mitad de las
veces, como al azar. La segunda es de diseño. Cuando el sitio servía el conjunto compra, la
variable \texttt{log\_num\_games\_owned} se llenaba con las compras al año que la persona declara
en el formulario, pero el modelo la había aprendido del tamaño de la biblioteca de quien escribió
la reseña; las compras declaradas por año eran un sustituto de ese dato que nunca se validó.

Con el conjunto juego, el riesgo es del título y vale igual para cualquier persona: no depende de
un dato que nadie verifica y se explica solo con propiedades del juego. El perfil que se declara
en el sitio sirve para la afinidad, que dice qué tanto encaja un juego con quien lo declaró, y no
mueve el riesgo.

\subsection{Validación agrupada por juego}\label{sec:modelacion-validacion}

Si las reseñas se repartieran al azar entre folds, el mismo modelo daría \cifraPRAUCKFold{}, con
una desviación de apenas \cifraPRAUCKFoldStd{}, porque cada fold de prueba tendría juegos que el
modelo ya vio (sección~\ref{sec:exploracion-juego}). Agrupando por \texttt{appid} da
\cifraPRAUCModelo{}~$\pm$~\cifraPRAUCModeloStd{}, la cifra que vale para un juego nuevo. El reparto
de los juegos también pesa: con \cifraParticionesAleatorias{} particiones por juego hechas al azar,
el PR-AUC medio va de \cifraPRAUCParticionMin{} a \cifraPRAUCParticionMax{}. Por eso la partición
está congelada en un archivo de \texttt{backend/referencias/}, y toda evaluación y los umbrales la
leen de ahí. Hay otra razón: GroupKFold reparte los juegos según su número de reseñas,
\cifraJuegosEmpatadosEnElTope{} de los \cifraJuegosEntrenamiento{} tienen exactamente el tope de
\cifraTopeIngesta{}, y cada versión de scikit-learn desempata distinto. Con la
\cifraVersionSklearnColab{}, la que trae Colab, solo \cifraJuegosMismoFoldColab{} juegos quedan en
el mismo fold; aun así, el PR-AUC es
\cifraPRAUCParticionAlternativa{}~$\pm$~\cifraPRAUCParticionAlternativaStd{} y el modelo le gana
al trivial en \cifraPliegosGanadosAlternativa{} de \cifraPliegues{} folds, igual que con la
congelada.

\subsection{Coeficientes y bandas de riesgo}\label{sec:modelacion-bootstrap}\label{sec:modelacion-umbrales}

Los coeficientes de la crítica están en la sección~\ref{sec:exploracion-critica}. Los demás cruzan
el cero: el del precio va de \cifraCoefPrecioICInf{} a \cifraCoefPrecioICSup{}, el del descuento,
de \cifraCoefDescuentoICInf{} a \cifraCoefDescuentoICSup{}, y el de la gratuidad, de
\cifraCoefGratisICInf{} a \cifraCoefGratisICSup{}. El del precio es positivo en el
\cifraCoefPrecioReplicasPositivas{} de las réplicas, pero con \cifraJuegosEntrenamiento{} juegos su
intervalo no descarta el cero, y por eso su evidencia se marca como débil.

El sitio no muestra el score: muestra una de tres bandas, bajo, medio o alto. Los cortes dividen en
tercios los scores fuera de pliegue\footnote{El score fuera de pliegue de una reseña lo da un modelo
entrenado sin el juego de esa reseña, con la misma partición congelada.} de las reseñas de data-v1:
un juego pasa a medio desde \cifraUmbralMedio{} y a alto desde \cifraUmbralAlto{}. Los tercios son
de reseñas, y cada juego se califica con el modelo entrenado con todo data-v1, así que los juegos no
se reparten en partes iguales: los \cifraJuegosEntrenamiento{} quedan en \cifraBajoEntrenamiento{},
\cifraMedioEntrenamiento{} y \cifraAltoEntrenamiento{}, y los \cifraJuegosCatalogo{} del catálogo,
en \cifraBajoCatalogo{}, \cifraMedioCatalogo{} y \cifraAltoCatalogo{}. Por los pesos balanceados, el
score no es una probabilidad calibrada: los cortes quedan muy por encima de la prevalencia, que es
de \cifraPrevalenciaEntrenamientoPorResena{}. Sirve para ordenar títulos, no como la probabilidad
de que una compra termine en la señal. Los umbrales salen solo de data-v1, y los
\cifraJuegosExternos{} externos se clasifican con ellos sin ajustarlos.

\subsection{Resultados por fold y prueba externa}\label{sec:evaluacion}\label{sec:evaluacion-folds}\label{sec:evaluacion-externa}

La figura~\ref{fig:pr-auc-por-fold} muestra el PR-AUC de cada fold de la partición congelada. El
modelo le gana al trivial en \cifraPliegosGanados{} de \cifraPliegues{} folds, pero por márgenes
muy distintos, de \cifraCocienteFoldMin{} a \cifraCocienteFoldMax{} veces; en promedio llega a
\cifraPRAUCModelo{} contra \cifraPRAUCTrivial{}, \cifraPRAUCCociente{} veces el trivial. Es una
desviación grande al lado de la media: el resultado depende de qué juegos le tocan a cada fold.

\begin{figure}[htbp]
  \centering
  \includegraphics[width=\textwidth]{pr-auc-por-fold.pdf}
  \caption{PR-AUC del modelo y del trivial en cada fold de la partición congelada (data-v1) y en
  la prueba externa. Sobre cada barra, cuántas veces supera el modelo al trivial. Fuente:
  \texttt{evaluar\_gkf} sobre data-v1 y \texttt{docs/evidencia/prueba-externa.json}.}
  \label{fig:pr-auc-por-fold}
\end{figure}

Los \cifraJuegosExternos{} títulos que llegaron con data-v2 se califican con el modelo entrenado en
data-v1, con sus umbrales y su mediana de Metacritic, igual que lo hace la API. Su PR-AUC es
\cifraPRAUCExterno{} contra \cifraPRAUCExternoTrivial{} del trivial: \cifraPRAUCExternoCociente{}
veces, la última pareja de la figura. Con los títulos remuestreados \cifraExternoReplicas{} veces,
el intervalo al \cifraNivelDeConfianza{} del cociente va de \cifraExternoICInf{} a
\cifraExternoICSup{}, y solo \cifraExternoReplicasSinVentaja{} réplicas quedan en 1 o menos. Hay
señal fuera de muestra, pero es modesta: el extremo inferior queda cerca de 1. Entrenar con los
\cifraJuegosCatalogo{} juntos bajaría el ruido entre folds, pero se perdería la única evaluación con
juegos que el modelo no vio, y por eso el modelo se queda entrenado con data-v1.

\subsection{Por qué cae fuera de muestra}\label{sec:evaluacion-caida}\label{sec:evaluacion-filo}

Una parte de la caída se ve en la variable que más pesa, la crítica: de los \cifraAltoCatalogo{}
juegos de riesgo alto del catálogo, \cifraAltoSinNotaCatalogo{} no tienen nota. Por reseña, en los
de entrenamiento la banda alta sin nota tiene \cifraTasaAltoSinNotaEntrenamientoPorResena{} contra
\cifraTasaBajoEntrenamientoPorResena{} de la baja, una tabla solo descriptiva porque el modelo ya
vio esos juegos. En los externos la distancia se acorta, \cifraTasaAltoSinNotaExternosPorResena{}
contra \cifraTasaBajoExternosPorResena{}, y la baja y la media se invierten
(\cifraTasaBajoExternosPorResena{} contra \cifraTasaMedioExternosPorResena{}). Con el promedio por
juego la lectura es la misma: \cifraTasaAltoSinNotaExternosPromJuegos{} contra
\cifraTasaBajoExternosPromJuegos{} en los externos. La falta de nota sigue señalando más riesgo,
pero separa menos en los títulos nuevos, que además traen otra mezcla
(sección~\ref{sec:exploracion-externos}; detalle en \texttt{docs/evidencia/metacritic-por-banda.md}).

Las bandas cortan un score continuo, y un juego cerca de un corte puede cambiar de banda con un
cambio mínimo en los datos o en el reentrenamiento. Los dos más cercanos son Hollow Knight, con
\cifraScoreHollowKnight{} contra el corte de \cifraUmbralMedio{}, en la banda baja, y Warframe,
con \cifraScoreWarframe{} contra \cifraUmbralAlto{}, en la alta. Por eso no se usan como casos de
demostración.

\subsection{Cómo se explica el riesgo}\label{sec:interpretabilidad}\label{sec:interpretabilidad-factores}\label{sec:interpretabilidad-avisos}

Para cada juego, la API calcula el aporte de cada variable al score, su coeficiente por su valor
estandarizado\footnote{El aporte está en log-odds, la escala en la que la regresión suma: positivo
sube el score y negativo lo baja.}, y la ficha muestra los tres de mayor aporte, aunque el primero
tenga evidencia débil. Cada factor lleva su nivel de evidencia, según sus intervalos
(sección~\ref{sec:modelacion-bootstrap}): sólida para la cobertura y la nota de la crítica, y débil
para la gratuidad, el precio y el descuento, con la leyenda «evidencia débil: con
\cifraJuegosEntrenamiento{} juegos no se distingue de cero». Junto a cada factor va la cifra del
juego y su referencia: la nota contra el promedio del catálogo, \cifraNotaPromedioCatalogo{}, y el
precio contra la mediana de los de pago, \cifraPrecioMedianoCatalogo{} pesos, que se parece más a
la escala logarítmica del modelo que el promedio.

Un factor con un aporte menor que \texttt{UMBRAL\_TIPICO}, \cifraUmbralTipico{}, se muestra sin
flecha, como «cerca de lo típico del catálogo». Ese valor es una decisión de producto, no una
propiedad estadística: no sale de un intervalo ni de una prueba. Se eligió para que ninguna flecha
contradiga la cifra que tiene al lado, porque el modelo compara cada nota contra su media de
entrenamiento, \cifraMediaNotaModelo{}, mientras que la ficha la compara contra el promedio del
catálogo; una nota entre las dos aparece «por encima del promedio» y aun así recibe un aporte
pequeño que sube el score. Con \cifraUmbralTipico{}, ningún factor con flecha cae en ese caso en
los \cifraJuegosCatalogo{} juegos, y el generador de este documento lo comprueba.

Dos casos llevan aviso. En un juego gratis, la gratuidad y el precio de 0 dicen lo mismo y tiran en
sentidos opuestos, así que se muestran como un solo factor, y el aviso agrega que en el
entrenamiento había solo \cifraGratisEntrenamiento{} juegos gratis y que, para ellos, el modelo
extrapola. En los \cifraPagoSinPrecioEntrenamiento{} juegos de pago sin precio, el aviso dice que
la estimación es menos confiable, porque el modelo tomó el precio como 0.

\subsection{Los motivos: una heurística de palabras clave}\label{sec:interpretabilidad-motivos}

\texttt{/explicacion} no usa un modelo. Busca palabras clave en inglés de \cifraCategoriasMotivos{}
categorías (rendimiento, bugs, dificultad, controles, contenido y precio) en las reseñas con la
señal de cada juego, y una reseña puede caer en varias o en ninguna. Si el juego tiene menos de
\cifraUmbralMinCasos{} reseñas con la señal, no muestra motivos, y cada porcentaje se calcula sobre
las reseñas que sí cayeron en alguna categoría, cuyo número la ficha dice. El límite es la
cobertura: de las \cifraNegativasTempranasLimpias{} negativas tempranas del conjunto limpio, el
\cifraCoberturaMotivosPorResena{} menciona al menos un motivo de la lista, y el
\cifraSinMotivoPorResena{} restante no menciona ninguno; rendimiento y bugs son los más frecuentes
(detalle en el notebook 00, §3.8). Por eso los motivos se leen como las quejas de las reseñas que
usan esas palabras, no como la causa de la señal en ese juego.

\subsection{Modelos de texto (Parte A)}\label{sec:modelacion-texto}

Los modelos que leen el texto de las reseñas, y con ellos la comparación entre algoritmos, se están
haciendo en la rama \texttt{modelos-texto}, con prerregistro\footnote{Un prerregistro fija, antes
de correr el análisis, qué se va a medir y qué resultado contará como éxito, para no acomodar la
pregunta a los resultados.}. Su resultado entrará a este documento cuando esa rama llegue a master.

\reservaPendiente{11cm}{resultado de la Parte A, de la rama \texttt{modelos-texto}}
